\subsection{Coherent coefficient kernels and their completion limits}
The coefficient-object alternative in \fid{F54} has an explicit
realization. It retains the actual rank-three word algebra on a
split Grothendieck ring, with natural maps for the presentation
relations. Its completion and adjoint behavior show why this
does not yet select the requested canonical basis.

\begin{finding}[\fid{F56}: coefficient kernels, split $K_0$, and completion]
\label{f:coefficientkernelcategory}
Let $\Lambda=\Z[v,v^{-1}]$ and use the actual positive
rank-three matrices of \fid{F43}.
There is a concrete additive monoidal category
$\mathcal P_{64}$ of finite graded matrix kernels, acting
on $64$ graded coefficient states, such that
\[
 K_0^{\mathrm{split}}(\mathcal P_{64})
   \simeq A_\Lambda^{\mathrm{act}}
   =\Lambda\langle P_1,P_2\rangle
       \subset M_{64}(\Lambda).
\]
It has explicit natural quadratic and degree-five
relation isomorphisms, with coherent associators and
ordered coefficient comparisons.
The algebra on the right is the actual word-algebra
image, not the full ambient matrix algebra.

The category is not closed under direct summands
or under the adjoints of its named generators.
Its Karoubi completion satisfies
\[
 K_0^{\mathrm{split}}\operatorname{Kar}(\mathcal P_{64})
       \simeq\Lambda^{24}\times M_2(\Lambda)^{20},
\]
of rank $104$ over $\Lambda$.
On the faithful six-state version, the completion instead has
\[
 \Lambda^2\times M_2(\Lambda)^2.
\]
Its generic rank is ten, but its localized order is the
split normalization of the actual colliding center.

For fixed $q$, replacing degree $r$ by a specified Weil
line of Frobenius $q^{r/2}$ gives an exact arithmetic
coefficient-object realization, preserving the natural
relation maps at every extension degree.
In odd characteristic, entrywise application of
\fid{F53}'s $J$ gives actual perverse convolution kernels
and maps, with the known weight of every homogeneous piece.
This is a rank-three graded construction, not a selection
of VIII's canonical simple labels or integral forms.
\end{finding}

\begin{proof}
\par\smallskip\noindent\textbf{The concrete matrix-kernel category.}\quad
Let $\mathcal V$ be finite-dimensional $\Z$-graded vector
spaces over a characteristic-zero field, with
degree-preserving maps and ordinary tensor symmetry.
There are no Koszul signs: the grading is formal
quantum degree.
Write $\Bbbk\langle d\rangle$ for a line in degree $d$,
and
\[
 \operatorname{ch}(U)=\sum_d\dim U_d\,v^d.
\]
For a finite state set $I$, define
$\operatorname{Mat}_I(\mathcal V)$ with entrywise maps and
convolution
\begin{equation}\label{eq:gradedmatrixconvolution}
 (K\star L)_{ij}=\bigoplus_h K_{ih}\otimes L_{hj}.
\end{equation}
Its unit is diagonal $\Bbbk$.
It acts on row vectors by
\[
 (T_K X)_j=\bigoplus_i X_i\otimes K_{ij},
 \qquad T_L\circ T_K\simeq T_{K\star L}.
\]

Use the $64$ actual \fid{F43} states: $20$ scalar-plus
states, $4$ scalar-minus states, $16$ copies of $N(\delta)$,
and $4$ copies of $N(D-2)$.
Set
\[
 A_1=\Bbbk\langle-1\rangle\oplus\Bbbk\langle1\rangle,\quad
 A_2=\Bbbk\langle-2\rangle\oplus\Bbbk\langle2\rangle,\quad
 \mathcal D=A_1\otimes A_1.
\]
Their characters are $\delta,D-2,D$.
On a two-state block use
\begin{equation}\label{eq:actualgradedkernels}
 K_1=\begin{pmatrix}\mathcal D&0\\A&0\end{pmatrix},
 \qquad
 K_2=\begin{pmatrix}0&A\\0&\mathcal D\end{pmatrix},
 \qquad A=A_1\ \hbox{or}\ A_2.
\end{equation}
These are the source's actual right-row matrices.
On plus states both kernels are $\mathcal D$; on minus
states they are zero.
The basis is a basis of the actual tensor lattice after
the proved finite localization.
No new equality with the original global tensor-monomial
lattice is asserted.

For any parenthesized convolution word, flattening gives
\[
 \bigoplus_{i_1,\ldots,i_{m-1}}
 (K_1)_{i_0i_1}\otimes\cdots\otimes(K_m)_{i_{m-1}i_m}.
\]
The associator preserves the state sequence and ordered
tensor factors.
Every route around its pentagon therefore has the same
action on every pure tensor in every summand.
The unit identities have the same verification.
Coefficient symmetries permute pure tensor positions.
Adjacent swaps satisfy their square, distant-commutation,
and braid relations because the resulting permutations agree.
This establishes the monoidal coherence used below.

\par\smallskip\noindent\textbf{Quadratic maps and the ordered overlap.}\quad
A two-state $K_i$ resets to state $i$, so
\[
 (K_i\star K_i)_{hi}=(K_i)_{hi}\otimes\mathcal D.
\]
The path identification and coefficient symmetry give
\begin{equation}\label{eq:kernelquadratic}
 K_i\star K_i\simeq K_i\otimes\mathcal D
                   \simeq\mathcal D\otimes K_i.
\end{equation}
These are actual natural maps.
The same description works on plus states and zero
objects on minus states.
For every $m\ge1$, retaining the initial edge and the
chronological list of loops gives
\[
 K_i^{\star m}\simeq
 K_i\otimes\mathcal D^{\otimes(m-1)}.
\]
Flattening makes this independent of convolution
parenthesization.

For precision, use the left-coefficient version of
\eqref{eq:kernelquadratic} and label a triple path
$u\otimes d_1\otimes d_2$.
Reducing the first pair and then the remaining pair
gives $d_1\otimes d_2\otimes u$.
Reducing the last pair, moving its coefficient left,
and then reducing the remaining pair gives
$d_2\otimes d_1\otimes u$.
The comparison is the transposition of the two
$\mathcal D$ factors.
It is nonidentity, witnessed by their distinct
degree-zero labels, but is precisely ordinary
symmetric bookkeeping, not a coherence obstruction.

In general a reduction tree gives a permutation
$\sigma_T$ of the labeled internal $\mathcal D$ factors.
Inductively, local maps only permute those labels and
retain the initial edge.
Compare two normal forms by
$\sigma_{T'}\sigma_T^{-1}$.
All comparison pentagons telescope to the identity,
and the elementary comparisons satisfy the symmetric-group
relations already checked.
This supplies coherent ordered comparisons; it does not
impose strict equality after discarding the labels.

\par\smallskip\noindent\textbf{The actual degree-five relation maps.}\quad
IV, Theorem~11.16, equations~(111)--(112), p.~65, gives
\[
 P_i^2=DP_i,\qquad
 P_1(P_2P_1-a_1^2)(P_2P_1-a_2^2)
 =
 P_2(P_1P_2-a_1^2)(P_1P_2-a_2^2),
\]
with $a_1=\delta$, $a_2=D-2$ \cite{IV}.
Let $V_1=K_1\star K_2\star K_1$ and
$W_1=K_1\star K_2\star K_1\star K_2\star K_1$;
use $V_2,W_2$ for their mirrored words.
Put
\[
 U=A_1^{\otimes2},\quad V=A_2^{\otimes2},\quad
 \mathcal S=U\oplus V,\quad \mathcal T=U\otimes V.
\]
The positive rearrangement has sides
\begin{align}
 L&=W_1\oplus(V_2\otimes\mathcal S)
                         \oplus(K_1\otimes\mathcal T),\nonumber\\
 R&=W_2\oplus(V_1\otimes\mathcal S)
                         \oplus(K_2\otimes\mathcal T).
 \label{eq:kerneldegreefive}
\end{align}
We specify an actual natural isomorphism.

On an $N(A)$ block, every internal alternating step
is the $A$ edge, giving chronological identifications
\[
 V_j\simeq K_j\otimes A^{\otimes2},\qquad
 W_j\simeq K_j\otimes A^{\otimes4}.
\]
Suppress tensor symbols in the following lists.
For $A=A_1$, expand the sides as
\[
 L=(K_1UU,K_2UU,K_2UV,K_1UV),\qquad
 R=(K_2UU,K_1UU,K_1UV,K_2UV).
\]
Map positions by $(1\,2)(3\,4)$, using identity
coefficient maps.
For $A=A_2$, the lists are
\[
 L=(K_1VV,K_2VU,K_2VV,K_1UV),\qquad
 R=(K_2VV,K_1VU,K_1VV,K_2UV).
\]
Send positions $1,2,3,4$ to $3,4,1,2$.
The first and third maps are identities; the other two
interchange the $U,V$ coefficient blocks.
On plus states use the identical expressions
$\mathcal D^5,\mathcal D^3\mathcal S,\mathcal D\mathcal T$,
and on minus states use the unique zero map.

These maps specify every summand and factor permutation.
They are degree preserving and natural.
Exchanging the generator labels gives the inverse:
the displayed permutations reverse, and the opposite
coefficient interchange inverts the first.
The same map on every multiplicity copy commutes with
its matrix units.
Tensoring with any kernel context gives natural relation
maps; disjoint applications commute by the tensor
interchange identity.
Associators and coefficient symmetries have the coherence
already proved.

Thus there is a coherent category/action with explicit
maps for both algebra presentation relations.
No source-prescribed additional higher $2$-relation
among overlapping degree-five maps has been assumed.
An assertion that every unspecified rewriting loop is
identity would be stronger and is not made here.

\par\smallskip\noindent\textbf{The split Grothendieck ring.}\quad
Define $\mathcal P_{64}$ to be the full additive
subcategory whose objects are finite direct sums of
grading shifts of convolution words in $I,K_1,K_2$,
including their isomorphic copies.
Do not close under arbitrary direct summands.
It is monoidal, and use its split exact structure.

Isomorphism classes of finite graded matrix kernels
are determined by entrywise graded dimension.
If those matrices agree, choose an isomorphism in every
homogeneous entry. Fullness puts the resulting
isomorphism in $\mathcal P_{64}$.
The object monoid therefore embeds in
$M_{64}(\mathbb N[v,v^{-1}])$ and is cancellative.
Its image is exactly the nonnegative Laurent
combinations of actual word matrices.

Group completion is injective into $M_{64}(\Lambda)$.
Indeed, equality of two formal differences gives,
after cross-addition, identical dimension matrices,
hence an actual kernel isomorphism and the same
split Grothendieck class.
The image is the $\Z$-span of all shifted words,
which is exactly $A_\Lambda^{\mathrm{act}}$.
Convolution respects multiplication, proving
\[
 K_0^{\mathrm{split}}(\mathcal P_{64})
          \simeq A_\Lambda^{\mathrm{act}}.
\]
All four actual generic simple types occur, so the
actual generic algebra acts faithfully and has
dimension ten.
Removing repeated copies does not change any word
relation. After the established finite localization,
this is the actual order of \fid{F19}.

The qualifier \emph{split} matters.
No abelian or semisimple simple-object decomposition,
idempotent completeness, or prescribed canonical basis
has been asserted for this category.

\par\smallskip\noindent\textbf{Karoubi completion on $64$ states.}\quad
The unit already has the independent state idempotents
in its endomorphism algebra.
They split in the Karoubi completion to give every
$E_{ii}$.
On a two-state block,
$E_{22}\star K_1\star E_{11}$ is the coefficient $A$
in position $(2,1)$.
Splitting a homogeneous basis line and shifting its
degree gives $E_{21}$.
Similarly $K_2$ gives $E_{12}$.
Together with the diagonal objects, these give every
elementary two-by-two kernel and its shifts.
A scalar state supplies its one elementary diagonal kernel.

Every original word is block diagonal with respect to
the $24$ scalar and $20$ two-state copies.
An entrywise idempotent image cannot create a missing
matrix entry.
Conversely every finite graded kernel on these blocks
is a sum of the elementary shifted kernels just constructed.
The ambient fullness of maps therefore gives
\[
 \operatorname{Kar}(\mathcal P_{64})
       \simeq\mathcal V^{24}\times
                    \operatorname{Mat}_2(\mathcal V)^{20}.
\]
Hence its split Grothendieck ring is
\[
 \Lambda^{24}\times M_2(\Lambda)^{20}.
\]
There are $104$ elementary simple-kernel families
up to shift and $64$ simple summands of the unit.
There are no new entries between distinct multiplicity
copies, so the result is not $M_{64}(\Lambda)$.
The resulting positive elementary basis is a basis
of this enlarged algebra, not of the actual word algebra.

\par\smallskip\noindent\textbf{The faithful six-state normalization.}\quad
Retain one copy of each scalar type and each two-state type.
Its word relations are exactly the same, so its
$\mathcal P_6$ has the same split Grothendieck algebra.
The same completion argument gives
\[
 \operatorname{Kar}(\mathcal P_6)
   \simeq\mathcal V^2\times\operatorname{Mat}_2(\mathcal V)^2,
 \qquad
 K_0^{\mathrm{split}}\operatorname{Kar}(\mathcal P_6)
   =\Lambda^2\times M_2(\Lambda)^2.
\]
Its rank is ten, but the relative order has changed.

Work over $R=\mathbb Q[v,v^{-1}]_{(v-1)}$.
The actual center maps to the split branch center by
\[
 R[x]/(x(x-\Delta))\longrightarrow R\times R,\qquad
 f_0+f_1x\longmapsto(f_0,f_0+\Delta f_1).
\]
Its image is precisely the pairs congruent modulo
$\Delta$, and its cokernel is $R/(\Delta)$.
The center cokernel has length two; the four matrix
entries give length eight.
After the regular branch basis changes, the completed
kernel algebra is the split normalization order
\[
 R^2\times M_2(R)^2.
\]
The actual special fiber has center
$\mathbb Q^2\times\mathbb Q[\epsilon]/(\epsilon^2)$
and a four-dimensional matrix radical.
The displayed map kills that radical on specialization
and has six-dimensional image.
The completed order has a reduced split center.
Thus equal generic rank does not retain the collision.

\par\smallskip\noindent\textbf{Adjoints require an enlargement.}\quad
Finite tensor-Hom adjunction gives both left and right
adjoint kernels
\[
 (K^\dagger)_{ij}=(K_{ji})^*.
\]
Graded dual reverses degree.
The chosen $\mathcal D,A_1,A_2$ are self-dual, so
\[
 K_1^\dagger=\begin{pmatrix}\mathcal D&A\\0&0\end{pmatrix},
 \qquad
 K_2^\dagger=\begin{pmatrix}0&0\\A&\mathcal D\end{pmatrix}.
\]
Their supports already show that the named generators
are not self-adjoint, even up to shift.

Neither adjoint belongs to $\mathcal P_{64}$ or
$\mathcal P_6$.
For a nonempty word $w=i_1\cdots i_L$, the rank-one
matrix multiplication of \fid{F43} gives
\[
 P_w=D^{L-1-t}a^t u_{i_1}e_{i_L}^{T},\qquad
 u_1=(D,a)^T,\quad u_2=(a,D)^T,
\]
where $t$ counts adjacent changes.
At $v=1$, $D=4$ and $a=2$.
Both rows of every nonempty word are nonzero; the
identity also has both rows nonzero.
Every nonzero object in the word category is a
positive sum of shifted words.
Its restriction to either two-state block therefore
has both rows nonzero, since shifts evaluate to one
and no cancellation is possible.
But $K_1^\dagger$ has zero second row, and
$K_2^\dagger$ zero first row.
This proves exclusion.

This statement concerns the positive word-object
category. It does not assert that the adjoint operator
lies outside every localization of the actual algebra.
The Karoubi completion contains the transpose-dual
kernels and their elementary summands.

\par\smallskip\noindent\textbf{Arithmetic coefficient objects.}\quad
Fix a prime power $q$, a prime $\ell$ different from its
characteristic, and the algebraic number
$\lambda=\sqrt q>0$ under a specified complex embedding.
Choose its $\ell$-adic realization in a finite
coefficient field $E$.
Let $L$ be the Weil line with geometric Frobenius $\lambda$.
For $E$-linear graded spaces, send
\[
 E\langle r\rangle\longmapsto L^{\otimes r}.
\]
This extends by sums and tensors to an exact symmetric
monoidal functor.
It is fully faithful onto the semisimple category of
these characters: $\lambda^r=\lambda^s$ forces $r=s$,
and Frobenius maps between distinct characters vanish.
The equal-degree maps are the original coefficient maps.

Every specified kernel map is degree preserving,
so it becomes a Frobenius morphism.
The lines are continuous unramified representations:
$\lambda$ is an $\ell$-adic unit, its powers lie in a
compact profinite unit group, and the homomorphism
from $\Z$ extends to its profinite completion.
For every $n$, in the chosen complex embedding,
\begin{equation}\label{eq:gradedcoefficienttrace}
 \Tr\left(F^n\mid\bigoplus_r V_r\otimes L^r\right)
       =\sum_r\dim(V_r)q^{rn/2}.
\end{equation}
Thus $\mathcal D$ becomes
$L^2\oplus E^{\oplus2}\oplus L^{-2}$, and the quantum
relation is a tensor relation with that coefficient object
at every extension degree.
It is not the fixed scalar $D(\sqrt q)$ on one finite space.

Sending degree one instead to $E(-1)$ gives a full
Tate step and the substitution $v=q^n$.
The half-step above is a Weil coefficient line;
it is not asserted to be a geometric Tate object.
Different degrees give mixed weights, and different
words have their own finite trace objects.
This is not the fixed-operator comparison excluded
by \fid{F42} and \fid{F54}.

\par\smallskip\noindent\textbf{Actual perverse coefficient kernels.}\quad
Now assume $q$ is odd.
Use \fid{F53}'s exact fully faithful symmetric strong
monoidal functor $J$ to its perverse Gaussian image,
with specified product $\star_\psi$ and image unit $J(E)$.
Its symmetry is the transported symmetry, including
the established correction to the naive odd-shift swap.

On matrices of image objects define
\[
 (\mathcal K\star\mathcal L)_{ij}
       =\bigoplus_h\mathcal K_{ih}
                      \star_\psi\mathcal L_{hj}.
\]
The unit is diagonal $J(E)$.
Entrywise $J$ after the arithmetic functor gives
\[
 J_I(K\star L)\simeq J_I(K)\star J_I(L).
\]
To check coherence, expand an iterated product by its
intermediate state sequence and use F53's strong-monoidal
comparison on each tensor summand.
Every associator pentagon agrees summand by summand;
direct sums preserve the equality.
Full faithfulness and exactness are entrywise.
The matrix action on the corresponding product of
perverse image categories follows from the same formula.

All constructed quadratic, degree-five, inverse,
tensor-context, and ordered-overlap maps are therefore
actual perverse morphisms with the same comparisons.
This uses the explicit coefficient maps, not an
inference from a Laurent $K_0$ equality.

A degree-$d$ line has Weil weight $d$ in the half-step
convention, and $J(L^d)$ is perverse pure of weight $d$.
In the full Tate convention its weight is $2d$.
An arbitrary coefficient is the specified finite sum
of these pieces; its total need not have one weight.
The split-$K_0$, completion and adjoint calculations
carry to this full image by equivalence.
They do not classify all perverse sheaves in the
ambient constructible category.
\end{proof}

\paragraph{The precise remaining categorical comparison.}
This construction gives actual coherent operations and
maps at rank three, including an all-extension arithmetic
coefficient-object interpretation.
It does not supply compatible positive graded kernels
through all tensor degrees, a source-prescribed higher
categorical presentation, or VIII's global/local lattice and canonical
simple-label comparisons.
The fixed-parameter all-degree realization of \fid{F53}
is a separate statement.

After forgetting arithmetic descent, the lines $L^d$
are ordinary constant coefficient lines and their Gaussian
images have the same underlying geometric rank-one
constituent.
Their degree labels remain arithmetic/grading data,
rather than a new family of geometric cell labels.
The completion and adjoint results show concretely
why the present positive word category is not already
the desired semisimple canonical-basis category.
They do not refute a different categorical construction.

% Separate paragraph for Part III.
\paragraph{\fid{F56} in Grothendieck's language.}
The effective object monoid, its split group completion,
and its Karoubi completion are distinct constructions.
Here cancellation and the exact algebra image are proved,
and the completion changes the positive generators and
the relative order.
SGA~6, V, \S\S2--3, with cote~133, scan~18, provides the
audited representation/$\lambda$/$\sigma$ vocabulary;
it applies here to the symmetric coefficient category.
No $\lambda$-ring structure is asserted on the
noncommutative split $K_0(\mathcal P_{64})$, and
no $\gamma$-filtration positivity is inferred
\cite{SGA6,Archive}.

The actual coherent algebra and its action use the
affine module framework of EGA~I, Theorem~1.3.7, p.~86,
Corollary~1.3.8, p.~87, and Theorem~1.4.1, pp.~90--92,
and the coaction language of SGA~3, I, 4.7.2--3,
with cote~90, scan~17 \cite{EGA,SGA3,Archive}.
The normalization formula and all length computations
have been proved explicitly.
The perverse realization uses the actual Gaussian
convolution and primary compact-trace formalism
already supplied in \fid{F53}.
None of these framework identifications makes the
elementary matrix-kernel basis into the designated
nonstandard canonical basis.

